\IfLanguageName{english}{%
    \selectlanguage{dutch}
    \chapter*{Samenvatting}
    \lipsum[1-4]
    \selectlanguage{english}
}{}
\chapter*{\IfLanguageName{dutch}{Samenvatting}{Abstract}}

Bezoekers van begraafplaatsen ervaren regelmatig moeilijkheden bij het
terugvinden van een specifiek graf, wat een stressvolle ervaring kan zijn in een
emotioneel gevoelige context. Bestaande digitale oplossingen zoals gemeentelijke
geoloketten bieden wel kaartnavigatie, maar vereisen dat de gebruiker zijn
aandacht afwisselt tussen scherm en omgeving. Deze bachelorproef onderzoekt of
augmented reality (AR) een efficiënter en intuïtiever alternatief kan bieden.

De centrale onderzoeksvraag luidt: \textit{Op welke manier kan digitalisatie,
    in combinatie met Augmented Reality, worden ingezet om de efficiëntie van het
    navigeren op een kerkhof te verbeteren?} Als antwoord werd een proof of concept
ontwikkeld: een native iOS-applicatie die een 3D AR-navigatiepijl en een
world-locked grafmarker combineert met een interactieve satellietminikaart.
De applicatie werd gebouwd met ARKit, SceneKit en CoreLocation, en geëvalueerd
op het kerkhof van Sint-Gillis-Waas aan de hand van drie testscenario's.

Naast de technische implementatie werd ook een juridische analyse uitgevoerd.
Deze toont aan dat grafdata juridisch veilig verwerkt kan worden voor
navigatiedoeleinden, mits het rijksregisternummer categorisch wordt uitgesloten.
Voor een productieomgeving zijn een verwerkersovereenkomst met de gemeente en
een DPIA aanbevolen.

De testresultaten tonen aan dat AR-navigatie een duidelijke meerwaarde biedt
ten opzichte van navigatie zonder digitale hulp of met enkel een kaart.
De applicatie werd positief beoordeeld door de testpersonen. Als
aanbevelingen voor verdere ontwikkeling worden een onboarding voor nieuwe
gebruikers, een duidelijkere pijlvorm en de integratie van OSM-paddata voor
een geroute navigatielijn naar voren geschoven.